\documentclass[preprint,11pt,authoryear]{elsarticle}

\usepackage{amsmath,amssymb,amsthm}
\usepackage{booktabs}
\usepackage{tikz}
\usepackage{pgfplots}
\pgfplotsset{compat=1.18}
\usepgfplotslibrary{groupplots}
\usepackage[margin=1in]{geometry}
\usepackage[hidelinks]{hyperref}
\hypersetup{pdftitle={Pricing the Reserve: Market Equilibrium, Lender Safety and Liquidity in an Uncollateralized On-Chain Lending Pool}, pdfauthor={Claude Code}}

\newtheorem{proposition}{Proposition}
\newtheorem{lemma}{Lemma}
\theoremstyle{definition}
\newtheorem{remark}{Remark}

\makeatletter
\newcommand{\inputrows}[1]{\@@input #1 }
\makeatother
\newcommand{\E}{\mathbb{E}}
\newcommand{\EL}{\mathrm{EL}}
\newcommand{\Prob}{\mathbb{P}}

% Generated by scripts/equilibrium.py; do not edit.
\newcommand{\Effr}{4.33}
\newcommand{\ARef}{12.33}
\newcommand{\ELre}{5.15}
\newcommand{\UBase}{85}
\newcommand{\PrimeIllus}{7.50}
\newcommand{\RC}{41.8}
\newcommand{\PremBps}{800}
\newcommand{\AFreeZero}{10.24}
\newcommand{\RBarZero}{50.3}
\newcommand{\AFreeOne}{11.42}
\newcommand{\RBarOne}{45.1}
\newcommand{\AFreeTwo}{12.60}
\newcommand{\RBarTwo}{40.9}
\newcommand{\IdleDrag}{0.76}
\newcommand{\MarginRoom}{2.09}
\newcommand{\RMaxLenders}{58.7}
\newcommand{\MarginAtFortyFive}{1.69}
\newcommand{\LenderMarginFortyFive}{1.43}
\newcommand{\ReserveTermFortyFive}{5.55}
\newcommand{\ReserveSurplusFortyFive}{0.40}
\newcommand{\MarginAtSixtyFive}{-0.78}
\newcommand{\LenderMarginSixtyFive}{-0.66}
\newcommand{\LRAtRC}{6.10}
\newcommand{\LRAtFortyFive}{5.76}
\newcommand{\GuarAtFortyFive}{5.76}
\newcommand{\GrowthAtFortyFive}{5.55}
\newcommand{\DoubleAtFortyFive}{12.5}
\newcommand{\PriceCreditAtFortyFive}{2.22}
\newcommand{\LRAtSixtyFive}{3.67}
\newcommand{\GuarAtSixtyFive}{3.67}
\newcommand{\GrowthAtSixtyFive}{8.01}
\newcommand{\DoubleAtSixtyFive}{8.6}
\newcommand{\PriceCreditAtSixtyFive}{1.54}
\newcommand{\GrowthGain}{2.5}
\newcommand{\SlopePerPoint}{0.10}
\newcommand{\GrowthContFortyFive}{5.71}
\newcommand{\GrowthDiscFortyFive}{5.69}
\newcommand{\GrowthContSixtyFive}{8.34}
\newcommand{\GrowthDiscSixtyFive}{8.32}
\newcommand{\TermsPerYear}{12.17}
\newcommand{\MVVarFortyFive}{0.15}
\newcommand{\IdentityErr}{4\times10^{-16}}
\newcommand{\PTenAtRC}{8.2}
\newcommand{\PTenAtFortyFive}{6.8}
\newcommand{\PTenAtFifty}{5.1}
\newcommand{\PTenAtFiftyFive}{3.9}
\newcommand{\PTenAtSixtyFive}{2.8}
\newcommand{\PTenAtFortyFiveSeedTwo}{4.2}
\newcommand{\PTenAtFortyFiveSeedFive}{2.0}
\newcommand{\PTenAtRCSeedFive}{2.9}
\newcommand{\MeanFortyFive}{5.46}
\newcommand{\MeanZero}{6.10}
\newcommand{\ResTenRC}{4.4}
\newcommand{\MeanRC}{5.67}
\newcommand{\ResTenFortyFive}{6.4}
\newcommand{\ResTenSixtyFive}{24.7}
\newcommand{\MVPenFive}{0.4}
\newcommand{\MVPenTwenty}{1.5}
\newcommand{\MVPenHundred}{7.7}
\newcommand{\PTenFortyFivePoolHundred}{14.1}
\newcommand{\PTenFortyFivePoolThousand}{7.5}
\newcommand{\PTenRCPoolHundred}{16.0}
\newcommand{\AnnuityAtE}{0.125}
\newcommand{\AnnuityAtFifteen}{0.199}
\newcommand{\SeedEquivFortyFive}{1.00}
\newcommand{\SeedCostFortyFive}{0.20}
\newcommand{\SeedCostEFortyFive}{0.13}
\newcommand{\ShareCostFortyFive}{0.34}
\newcommand{\ShareCostTenFortyFive}{0.21}
\newcommand{\SeedEquivFiftyFive}{3.75}
\newcommand{\SeedCostFiftyFive}{0.75}
\newcommand{\SeedCostEFiftyFive}{0.47}
\newcommand{\ShareCostFiftyFive}{1.39}
\newcommand{\ShareCostTenFiftyFive}{1.05}
\newcommand{\LateShare}{10}
\newcommand{\LateCollectedByTerm}{0.915}
\newcommand{\LateRemaining}{0.085}
\newcommand{\VFixedFortyFiveYield}{0.917}
\newcommand{\VFixedFortyFiveBase}{1.000}
\newcommand{\VFixedSixtyFiveYield}{0.000}
\newcommand{\VFixedSixtyFiveBase}{0.000}
\newcommand{\RHatBase}{46.0}
\newcommand{\VMaxBase}{1.006}
\newcommand{\VFortyFiveBase}{1.006}
\newcommand{\VSixtyFiveBase}{0.784}
\newcommand{\AFortyFiveBase}{12.26}
\newcommand{\ASixtyFiveBase}{15.74}
\newcommand{\VFortyFiveYield}{0.968}
\newcommand{\VFiftyYield}{0.912}
\newcommand{\VSixtyFiveYield}{0.704}
\newcommand{\AFortyFiveYield}{12.74}
\newcommand{\ASixtyFiveYield}{17.51}
\newcommand{\RHatMin}{42.2}
\newcommand{\RHatMax}{49.0}
\newcommand{\SeedGainTwo}{3.5}
\newcommand{\SeedGainFive}{7.6}
\newcommand{\RHatSeedTwo}{45.2}
\newcommand{\RHatSeedFive}{45.2}
\newcommand{\RHatAbarFive}{49.0}
\newcommand{\RHatAbarTwenty}{42.2}
\newcommand{\LoNinetyNine}{41.2}
\newcommand{\HiNinetyNine}{50.2}
\newcommand{\RHatHTen}{45.8}
\newcommand{\VMaxHTen}{1.005}
\newcommand{\LoNinetyNineHTen}{40.2}
\newcommand{\HiNinetyNineHTen}{50.2}
\newcommand{\VZeroHTen}{0.936}
\newcommand{\VFortyFiveHTen}{1.005}
\newcommand{\VSixtyFiveHTen}{0.789}
\newcommand{\AFortyFiveHTen}{12.27}
\newcommand{\RHatHTenYield}{0.0}
\newcommand{\VZeroHTenYield}{1.033}
\newcommand{\VFortyFiveHTenYield}{0.980}
\newcommand{\VSixtyFiveHTenYield}{0.720}
\newcommand{\RHatHFive}{39.0}
\newcommand{\LoNinetyNineHFive}{29.0}
\newcommand{\HiNinetyNineHFive}{46.2}
\newcommand{\RHatHTenKTwo}{45.2}
\newcommand{\RHatHMin}{39.0}
\newcommand{\RHatHMax}{48.8}
\newcommand{\YTenAtRC}{5.67}
\newcommand{\YTenAtZero}{6.10}
\newcommand{\YTenAtFortyFive}{5.46}
\newcommand{\YLongAtRC}{6.10}
\newcommand{\RHatSpread}{0.3}
\newcommand{\PlateauLoSpread}{0.3}
\newcommand{\PlateauHiSpread}{0.0}
\newcommand{\VolPerPointFifty}{1.3}
\newcommand{\IncidenceErr}{0.00}
\newcommand{\WaitFull}{30}
\newcommand{\WaitFullCushion}{26}
\newcommand{\WaitHalf}{12.4}
\newcommand{\WaitQuarter}{3.5}
\newcommand{\WaitThreeQuarter}{21.2}
\newcommand{\WaitHalfCap}{13.3}
\newcommand{\TermDays}{30}
\newcommand{\UMaxHalfWeek}{65}
\newcommand{\UMaxHalfWeekCushion}{78}
\newcommand{\AFreeAtULiq}{11.79}
\newcommand{\AFreeAtCap}{9.96}
\newcommand{\KappaLoss}{1.05}
\newcommand{\GammaBackward}{0.95}
\newcommand{\GammaProtectFortyFive}{0.43}
\newcommand{\GammaProtectSixtyFive}{0.62}


\journal{Journal of Financial Intermediation}

\begin{document}

\begin{frontmatter}

\title{Pricing the Reserve: Market Equilibrium, Lender Safety and Liquidity\\in an Uncollateralized On-Chain Lending Pool}

\author{Claude Code\fnref{fn1}}
\fntext[fn1]{Claude Code is an AI agent made by Anthropic, working with the Microcredit Protocol
Project; the research was directed by the project's human operator. Sources, data and code:
\url{https://github.com/scottonchain/microcredit-theory}.}
\affiliation{organization={Microcredit Protocol Project}, country={International}}

\begin{abstract}
A lending pool on a public blockchain lends stablecoins without collateral and sets aside a share of
every interest payment in a first-loss reserve that it never releases. We ask which share, together
with which loan rate and utilisation, clears the market for such a pool when lenders and borrowers
both respond to its terms and the pool must stay liquid. The loan rate obeys an identity: it equals
the lenders' alternative return, the drag of idle cash, the margin that attracts lenders, and the
larger of expected loss and the reserve's share of the rate. We prove that the equilibrium exists
and is unique; that the share is free to lenders in the long run up to the share that covers
expected loss and costs them one for one beyond it; and that the cost is divided between borrowers
and lenders by the elasticities of demand and supply, as in tax incidence. When lenders also care
about the chance of losing principal, a volume-maximising share exists that is never below the share
that covers expected loss, and a first-order condition characterises it. For a book with a 5\%
annual default probability and an 800 basis point premium, loan volume stays within 1\% of its
maximum for shares of \LoNinetyNine\ to \HiNinetyNine\%, a seed of outside capital of 1\% of deposits
buys the protection of three points of share at less than half the cost, and the pool's own
structure bounds the time to exit at one loan term. The elasticities are scenarios, not estimates.
\end{abstract}

\begin{keyword}
financial intermediation \sep loan pricing \sep first-loss capital \sep tax incidence \sep liquidity \sep decentralized finance
\end{keyword}

\end{frontmatter}

% ---------------------------------------------------------------------------------------------
\section{Introduction}

A lending pool on a public blockchain takes deposits in a stablecoin and lends them for 30 days
without collateral. Its borrowers lack a credit history, stable banking or existing digital assets,
and the pool lends against credit that an issuer grants, that others stake or back, or that
borrowers earn by paying interest \citep{conservedcredit2026}. To protect lenders, the pool sets
aside a share $r$ of every interest payment in a first-loss reserve, which pays losses before
lenders do. The protocol never releases a reserve that interest has funded, because the same share
is credited to the borrower as credit it can draw on, and releasing it would let that credit be
recaptured. How large should $r$ be?

The question is not only one of risk. A larger share protects lenders but costs them: interest
that goes to the reserve is interest they do not receive, and a reserve that is never released
grows without limit once its inflow exceeds losses \citep{pricingreserve2026}. Lenders can move
deposits to other pools. If the net yield falls, deposits fall, the pool lends less, and borrowers
who have nowhere else to borrow are left without a loan. The loan rate can rise to compensate, but
borrowers respond to the rate. And the pool has to stay liquid: lenders can withdraw on demand while
borrowers hold their loans for a month. This paper puts these forces in one model and asks which
share, loan rate and utilisation clear the market and meet a liquidity target.

A natural first guess is that the loan rate is a floor, such as the prime rate, plus a charge for the
reserve, plus a margin that attracts lenders. Proposition~\ref{prop:identity} makes that exact, with
two corrections. The floor that matters is the lenders' alternative return, not a bank's lending rate,
and every term is scaled by utilisation and by the share of interest that lenders keep.

\paragraph{Results} We prove five groups of results and compute them for a scenario.
\begin{enumerate}
\item \emph{The rate identity} (Proposition~\ref{prop:identity}). Any loan rate $a$ at which lenders
  earn $e+x$ satisfies $(1-f)\,a = (e+x)/u + \max\{\EL,\ a\,r\}$, with funding rate $e$, fee share
  $f$, utilisation $u$ and expected loss $\EL$. The reserve share costs lenders nothing in the long
  run while $a\,r\le\EL$ and costs them $u\,a$ per unit of share beyond.
\item \emph{Equilibrium and incidence} (Propositions~\ref{prop:eq} and~\ref{prop:incidence}). The
  market has a unique equilibrium. A rise in $r$ raises the equilibrium rate and lowers volume, and
  borrowers bear the share $\varepsilon/(\eta+\varepsilon)$ of the cost through the rate, with $\eta$
  and $\varepsilon$ the elasticities of demand and supply.
\item \emph{Lender safety and the volume-maximising share} (Lemma~\ref{lem:mono} and
  Proposition~\ref{prop:rhat}). If lenders also require a low probability of losing principal, some
  volume-maximising share is never below the share that covers expected loss, and a seed of outside
  first-loss capital raises volume at every share.
\item \emph{Liquidity} (Proposition~\ref{prop:liquidity}). A run on a pool of 30-day loans is paid
  within one loan term, and the utilisation that meets a target for the time to exit is explicit.
\item \emph{Earned credit} (Proposition~\ref{prop:growth}). A borrower's credit from dues grows at
  the rate $r\,a$ a year, so the reserve share is a slow engine of borrower growth.
\end{enumerate}
In a scenario with a 5\% annual default probability and the production premium of
\PremBps\ basis points (a loan rate of \ARef\%), the share that covers expected loss is \RC\%,
loan volume stays within 1\% of its maximum for shares between \LoNinetyNine\% and \HiNinetyNine\%,
and at 65\% it has fallen to \VSixtyFiveBase\ of its level at $r_c$ in the base case
(Section~\ref{sec:numbers}). A seed of 1\% of deposits replaces three points of share at less than half the
cost. The best share moves with the book's default risk: it is about four to six points above the share that
covers expected loss in each of three books.

\paragraph{Related work} Raising a capital buffer raises the price of credit when equity is dearer
than funding \citep{bcbs2010}. We obtain the same logic for a buffer that lenders fund from their
own interest and cannot withdraw. The incidence of that cost among lenders and borrowers follows
the elasticity formula of tax incidence \citep{harberger1962,fullerton2002}, with the reserve share
as the wedge between the loan rate and the lenders' yield. When higher rates raise default, lenders'
returns can bend back and credit is rationed \citep{stiglitz1981}. We treat that case as an
extension (Section~\ref{sec:limits}), because the reserve shields lenders from it as long as the
reserve covers losses. Evidence on the elasticity of loan demand among borrowers with little access
to credit shows demand curves that slope down, steeply above a lender's usual rate
\citep{karlan2008}, with a mean elasticity near $-1$ among clients of one microfinance institution
\citep{bogan2015}. On the lenders' side, we model a probability of losing principal that lenders will
tolerate, in the safety-first tradition \citep{roy1952,telser1955}, and a requirement that exit
be possible in a stated time, which echoes the liquidity constraints of intermediaries that fund
illiquid loans with demandable claims \citep{diamond1983,holmstrom1998}. Decentralized lending pools
set rates by utilisation \citep{leshner2019compound,aave2020,gudgeon2020}; this pool fixes the rate
at origination and limits utilisation by a cap instead.

The next section lists the features of the pool that the model takes as given. Sections~\ref{sec:model} to~\ref{sec:growth} build the model and prove the results, and
Section~\ref{sec:numbers} computes them. Section~\ref{sec:limits} states the limits, and
Section~\ref{sec:implications} draws the implications for the choice of the share.

% ---------------------------------------------------------------------------------------------
\section{The pool}\label{sec:pool}

The model uses the following features of the protocol's contract \citep{microcredit2026contract}.
\begin{itemize}
\item \emph{Loans and the rate.} Loans last $T=30$ days, accrue simple interest, and carry a fixed annual
  rate $a=e+s$ that the pool sets at origination: the funding rate $e$ (the effective federal funds
  rate, 4.33\% when deployed) plus a premium $s$. A share $f$ of interest goes to the protocol.
\item \emph{The reserve.} A share $r\in[0,1-f)$ of each interest payment joins the first-loss reserve $R$,
  which pays default losses before lenders do. Interest-funded reserve is never released. Anyone may
  add capital with \texttt{fundReserve}, which we call a seed. The contract would release seed
  capital only above all interest-funded reserve and provisions, so we treat it as locked.
\item \emph{Earned credit.} The share $r$ of interest that a borrower pays is added to its dues, which
  count as credit the borrower can draw on.
\item \emph{Utilisation.} New loans must keep lent principal at or below 90\% of lenders' assets and must
  leave 5\% of assets in cash. Lenders' assets exclude the reserve, which is a junior claim.
\item \emph{Liquidity.} The reserve's cash stays in the pool's cash balance. It therefore counts toward the
  buffer for new loans and toward cash available for withdrawals, but not toward the base on which the
  utilisation cap is computed. Withdrawals that cannot be paid at once join a queue, and a new loan
  must leave enough cash for everything queued, so repayments reach the queue first.
\end{itemize}
Table~\ref{tab:inputs} lists the inputs of the numerical scenario. They are scenarios, not estimates:
the pool has no default history, and the elasticities below have not been measured for it.

\begin{table}[ht]
  \centering\small
  \caption{Inputs of the scenario that the numerical sections use.}
  \label{tab:inputs}
  \begin{tabular}{p{0.34\linewidth}p{0.58\linewidth}}
    \toprule
    Input & Value \\
    \midrule
    Funding rate $e$, fee $f$ & \Effr\%, $0$ \\
    Premium and loan rate $a$ & \PremBps\ bps, \ARef\% \\
    Utilisation $u$ & \UBase\% (the contract's cap is 90\%) \\
    Default risk & annual probability 5\%, one-factor model at the Basel retail correlation, re-lent;
      expected loss $\EL=\ELre\%$ per exposure-year \citep{pricingreserve2026} \\
    Share that covers expected loss & $r_c=\EL/a=\RC\%$ \\
    Safety horizon $H$ & 10 years \\
    Borrowers' demand & constant elasticity $\eta=1$ (varied) \\
    Lenders' margin $\theta$ & exponential with mean $\tau=1\%$ (varied) \\
    Tolerated chance $\alpha$ of losing principal & exponential with mean $\bar\alpha=10\%$ (varied) \\
    \bottomrule
  \end{tabular}
\end{table}

% ---------------------------------------------------------------------------------------------
\section{The model}\label{sec:model}

\subsection{Lenders' return}

Per unit of deposits, the pool lends a fraction $u$ and collects interest $u\,a$ a year. In year $t$
a loss rate $L_t\in[0,\bar L]$ on lent principal arises, with mean $\E L_t=\EL$. The reserve $R_t$ per
unit of deposits starts at the seed $k\ge 0$, gains $u\,a\,r$ a year, and pays $u\,L_t$ first. Lenders
receive interest $u\,a\,(1-f-r)$ less the loss the reserve cannot pay:
\begin{equation}\label{eq:recursion}
  R_t=\max\{R_{t-1}+u\,a\,r-u\,L_t,\,0\},\qquad \ell_t=\max\{u\,L_t-R_{t-1}-u\,a\,r,\,0\},\qquad Y_t=u\,a\,(1-f-r)-\ell_t,
\end{equation}
where $Y_t$ is lenders' return in year $t$ per unit of deposits.

\begin{lemma}[Lenders' return]\label{lem:return}
For every path, $\sum_{t\le T}Y_t=T\,u\,a\,(1-f)-u\sum_{t\le T}L_t-(R_T-k)$. If the $L_t$ are independent
with common mean $\EL$, then almost surely
\begin{equation}\label{eq:longrun}
  \frac1T\sum_{t\le T}Y_t\longrightarrow y(a,r):=u\bigl[a\,(1-f)-\max\{\EL,\ a\,r\}\bigr].
\end{equation}
\end{lemma}

The first statement is an accounting identity: lenders' cumulative return is interest net of the fee,
less losses, less what the reserve has accumulated. The second says that, in the long run, a share
$r\le\EL/a$ costs lenders nothing, because the reserve is emptied by losses as fast as it fills, and a
share above that costs them $u\,a$ per unit of share, because the excess accumulates in a balance
they cannot receive. Lemma~\ref{lem:return} restates and proves, for independent years, the
long-run result of \citet{pricingreserve2026}; the proof is in the appendix.

\subsection{Participation}

Lenders differ in what they require. A lender of type $(\theta,\alpha)$ deposits if the long-run
yield covers the funding rate plus its required margin, $y(a,r)\ge e+\theta$, and if the probability
$p_H(r,k,a)$ that its return is negative in some year of the first $H$ years does not exceed its
tolerance, $p_H\le\alpha$. Types are independent. Let $G$ be the distribution of $\theta$ and
$A(p)=\Prob(\alpha\ge p)$ the share of lenders (by capacity $\bar V$) who tolerate a probability $p$,
with $A$ nonincreasing and $A(0)=1$. The volume that lenders fund is
\begin{equation}\label{eq:supply}
  V^s(y,p)=\bar V\,G(y-e)\,A(p),
\end{equation}
that is, utilisation times deposits. Without a safety requirement, $A\equiv1$, and the volume
$V^s(y)=\bar V\,G(y-e)$ has the inverse $\psi(V)=G^{-1}(V/\bar V)$: the margin over the funding rate
that the marginal lender requires to fund volume $V$. The elasticity of supply is
$\varepsilon=d\ln V^s/d\ln y=y/(V\psi'(V))$.

\subsection{Borrowers}

Borrowers' demand $V^d(a)$ for loans is continuous and strictly decreasing in the loan rate, with
elasticity $\eta=-d\ln V^d/d\ln a$. One study of the clients of a microfinance institution puts the mean
elasticity near $-1$ \citep{bogan2015}, and a field experiment found demand curves that slope downward, steeply for
increases above the lender's usual rate \citep{karlan2008}. Neither is a measurement for this pool.

\subsection{Liquidity}

Lenders can withdraw on demand. A run of a fraction $W$ of deposits at one moment is paid first
from cash, $(1-u)+c$ per unit of deposits, where $c=R/D$ is the reserve's cash cushion, and then
from repayments as loans mature. Section~\ref{sec:liquidity} derives the time to exit.

\subsection{Equilibrium}

Given the pool's choices of the reserve share $r$, seed $k$ and utilisation $u$, an
\emph{equilibrium} is a loan rate $a$ and a volume $V$ with $V=V^d(a)$ and $V=V^s(y(a,r),p_H(r,k,a))$:
lenders fund exactly what borrowers demand at the rate. If the pool sets a rate other than the
equilibrium rate, volume is the smaller of supply and demand.

% ---------------------------------------------------------------------------------------------
\section{The rate and its floor}\label{sec:floor}

\begin{proposition}[Rate identity and floor]\label{prop:identity}
Let lenders earn $e+x$ in the long run, $y(a,r)=e+x$ with $x\ge0$. Then
\begin{equation}\label{eq:identity}
  (1-f)\,a\;=\;\frac{e}{u}\;+\;\frac{x}{u}\;+\;\max\{\EL,\ a\,r\}.
\end{equation}
The lowest loan rate that gives lenders $e+x$ is
\begin{equation}\label{eq:floor}
  \underline a(r;x)=\max\Bigl\{\frac{(e+x)/u+\EL}{1-f},\ \frac{e+x}{u\,(1-f-r)}\Bigr\}.
\end{equation}
It is constant in $r$ up to $\bar r(x)=\EL/a_0(x)$, where $a_0(x)=((e+x)/u+\EL)/(1-f)$, and beyond it
is strictly increasing and convex, with slope $\underline a/(1-f-r)$, and tends to infinity as
$r\to 1-f$.
\end{proposition}

Read \eqref{eq:identity} as the sentence ``the rate is the floor, the idle-cash drag, the margin that
attracts lenders, and the reserve''. The floor is $e$, lenders' alternative. Because only a fraction $u$
of deposits earns interest, lenders need $e/u$ per dollar lent to earn $e$ per dollar deposited, so the
drag is $e(1/u-1)$ and the margin is $x/u$. The last term is the interest withheld from lenders: expected
loss, or the reserve share of the rate if that is larger. With $f=0$, at the production rate and $r=45\%$,
the terms are $\Effr\%$, $\IdleDrag\%$, $\MarginAtFortyFive\%$ and $\ReserveTermFortyFive\%$, and at
$r=65\%$ the margin is $\MarginAtSixtyFive\%$: lenders earn less than the funding rate
(Table~\ref{tab:decomposition}). The share at which lenders earn exactly $e$ at the production rate is
$1-f-e/(u\,a)=\RMaxLenders\%$.

\begin{table}[ht]
  \centering\small
  \caption{Decomposition of the production loan rate, $f=0$ (identity~\eqref{eq:identity}). Columns: reserve
  share $r$ (\%); funding rate $e$; idle-cash drag $e(1/u-1)$; lenders' margin per dollar lent, $x/u$;
  interest withheld, $\max\{\EL,a\,r\}$, of which above expected loss; lenders' long-run yield $e+x$ (\%). The terms
  of the first four sum to $\ARef\%$.}
  \label{tab:decomposition}
  \begin{tabular}{rrrrrrr}
    \toprule
    $r$ & $e$ & drag & $x/u$ & withheld & above $\EL$ & yield \\
    \midrule
    \inputrows{data/table_decomposition.tex}
    \bottomrule
  \end{tabular}
\end{table}

Figure~\ref{fig:floor} and Table~\ref{tab:floor} give the floor $\underline a(r;h)$ for lenders who require a
margin $h$ over the funding rate. At no margin it is flat at \AFreeZero\% up to $r=\RBarZero\%$ and rises after.
The production rate leaves \MarginRoom\ points of room above it, which is what a reserve share between $\RC\%$ and
$\RMaxLenders\%$ spends. For scale, the prime rate was \PrimeIllus\% when the funding rate was
4.33\% (an external input: by convention, the upper bound of the federal funds target range plus three points); the pool's floor lies
above it, because expected loss alone is $\ELre\%$.

\begin{figure}[ht]
  \centering
  % Figure: lowest APR at which lenders earn the funding rate plus a margin h, against the reserve share
% (data/floor_curve.csv).
\begin{tikzpicture}
  \begin{axis}[
      width=0.82\linewidth, height=6.2cm,
      xmin=0, xmax=0.7, ymin=7, ymax=26,
      xlabel={reserve share $r$}, ylabel={loan APR (\%)},
      tick label style={font=\footnotesize}, label style={font=\small},
      legend style={font=\scriptsize, draw=none, fill=white, fill opacity=0.85, text opacity=1, at={(0.03,0.97)}, anchor=north west},
    ]
    \addplot[thick, green!45!black] table[x=r, y=h0, col sep=comma] {data/floor_curve.csv};
    \addlegendentry{lenders' margin $h=0$}
    \addplot[thick, blue!65!black] table[x=r, y=h1, col sep=comma] {data/floor_curve.csv};
    \addlegendentry{$h=1\%$}
    \addplot[thick, red!70!black] table[x=r, y=h2, col sep=comma] {data/floor_curve.csv};
    \addlegendentry{$h=2\%$}
    \addplot[dashed, black, domain=0:0.7] {12.33};
    \node[font=\scriptsize, anchor=south east] at (axis cs:0.69,12.45) {production APR};
    \addplot[dotted, black, domain=0:0.7] {7.5};
    \node[font=\scriptsize, anchor=south east] at (axis cs:0.69,7.6) {prime rate};
  \end{axis}
\end{tikzpicture}

  \caption{The lowest APR at which lenders earn the funding rate plus a margin $h$, against the reserve
  share (Proposition~\ref{prop:identity}); $u=\UBase\%$, $\EL=\ELre\%$.}
  \label{fig:floor}
\end{figure}

\begin{table}[ht]
  \centering\small
  \caption{The floor $\underline a(r;h)$ (\%) by reserve share $r$ (\%) and lenders' margin $h$.}
  \label{tab:floor}
  \begin{tabular}{rrrr}
    \toprule
    $r$ & $h=0$ & $h=1\%$ & $h=2\%$ \\
    \midrule
    \inputrows{data/table_floor.tex}
    \bottomrule
  \end{tabular}
\end{table}

% ---------------------------------------------------------------------------------------------
\section{Equilibrium and incidence}\label{sec:eq}

\begin{proposition}[Equilibrium]\label{prop:eq}
Without a safety requirement ($A\equiv1$), let $\psi$ be continuous and strictly increasing with $\psi(0)=0$
and $\psi(V)\to\infty$ as $V\to\bar V$, and let $V^d$ be continuous and strictly decreasing with $V^d(a)\to0$
as $a\to\infty$. For every
$r\in[0,1-f)$ there is exactly one equilibrium $(a^*,V^*)$. It satisfies
\begin{equation}\label{eq:eqrate}
  (1-f)\,a^*=\frac{e+\psi(V^*)}{u}+\max\{\EL,\ a^*r\},\qquad V^*=V^d(a^*).
\end{equation}
The equilibrium rate $a^*$ is nondecreasing in $r$, and strictly increasing wherever $a^*r>\EL$. Volume
$V^*$ and lenders' margin $\psi(V^*)$ are nonincreasing in $r$.
\end{proposition}

\begin{proposition}[Incidence]\label{prop:incidence}
Where $a^*r>\EL$ and the primitives are differentiable,
\begin{equation}\label{eq:incidence}
  \frac{d\ln a^*}{dr}=\frac{\varepsilon}{\eta+\varepsilon}\cdot\frac1{1-f-r},\qquad
  \frac{d\ln V^*}{dr}=-\frac{\eta\,\varepsilon}{\eta+\varepsilon}\cdot\frac1{1-f-r},\qquad
  \frac{d\ln y^*}{dr}=-\frac{\eta}{\eta+\varepsilon}\cdot\frac1{1-f-r},
\end{equation}
where $\eta$ and $\varepsilon$ are the elasticities of demand and supply at the equilibrium and $y^*=y(a^*,r)$
is lenders' yield.
\end{proposition}

The reserve share acts as a tax of rate $r/(1-f)$ on lenders' interest, collected into a balance they cannot
draw. As with any tax, its burden is divided by elasticities. If lenders are perfectly elastic ($\varepsilon\to\infty$),
borrowers bear all of it: the rate rises by $\underline a/(1-f-r)$ per unit of share, the floor of
Proposition~\ref{prop:identity}. If borrowers are perfectly elastic, lenders bear all of it. Figure~\ref{fig:market} shows the
mechanism: the share lifts the supply curve once the reserve term in \eqref{eq:identity} exceeds expected loss, and the
equilibrium moves up the demand curve to a higher rate and a smaller volume. In the numerical scenario the formula
\eqref{eq:incidence} matches the numerical derivative of the equilibrium rate to two decimals of a percent
(Table~\ref{tab:incidence}).

\begin{figure}[ht]
  \centering
  % Figure: borrowers' demand and lenders' supply of loans, by reserve share (data/supply_demand.csv).
\begin{tikzpicture}
  \begin{axis}[
      width=0.82\linewidth, height=6.4cm,
      xmin=0.55, xmax=1.2, ymin=8, ymax=22,
      xlabel={loan volume (production volume at $r_c$ $=1$)}, ylabel={loan APR (\%)},
      tick label style={font=\footnotesize}, label style={font=\small},
      legend style={font=\scriptsize, draw=none, fill=white, fill opacity=0.85, text opacity=1, at={(0.03,0.97)}, anchor=north west},
    ]
    \addplot[thick, black] table[x=v, y=demand, col sep=comma] {data/supply_demand.csv};
    \addlegendentry{borrowers' demand}
    \addplot[thick, green!45!black] table[x=v, y=s_rc, col sep=comma] {data/supply_demand.csv};
    \addlegendentry{lenders' supply, $r=41.8\%$}
    \addplot[thick, blue!65!black] table[x=v, y=s_r50, col sep=comma] {data/supply_demand.csv};
    \addlegendentry{$r=50\%$}
    \addplot[thick, red!70!black] table[x=v, y=s_r60, col sep=comma] {data/supply_demand.csv};
    \addlegendentry{$r=60\%$}
    \addplot[only marks, mark=*, mark size=1.6pt, black, forget plot] coordinates {(1,12.33)};
  \end{axis}
\end{tikzpicture}

  \caption{Borrowers' demand (constant elasticity $\eta=1$) and lenders' supply of loans for three reserve shares,
  without a safety requirement ($\tau=1\%$). The dot is the production equilibrium at $r_c$.}
  \label{fig:market}
\end{figure}

\begin{table}[ht]
  \centering\small
  \caption{Incidence check, $\eta=1$, $\tau=1\%$. Columns: reserve share (\%); supply elasticity $\varepsilon$ at
  the equilibrium; borrowers' share of the cost through the rate, $\varepsilon/(\eta+\varepsilon)$ (\%);
  $d\ln a^*/dr$ by numerical differentiation; the formula \eqref{eq:incidence}.}
  \label{tab:incidence}
  \begin{tabular}{rrrrr}
    \toprule
    $r$ & $\varepsilon$ & borrowers' share & numerical & formula \\
    \midrule
    \inputrows{data/table_incidence.tex}
    \bottomrule
  \end{tabular}
\end{table}

% ---------------------------------------------------------------------------------------------
\section{Lender safety and the volume-maximising share}\label{sec:safety}

Expected return does not capture what a lender who thinks of its deposit as cash will accept. We use a
safety-first criterion: lender $\alpha$ participates if the probability that its return is negative in some
year of the first $H$ years does not exceed $\alpha$.

\begin{lemma}[Monotone safety]\label{lem:mono}
The principal-loss event $\{Y_t<0\text{ for some }t\le H\}$ equals $\{u\,L_t>u\,a\,(1-f)+R_{t-1}\text{ for some }t\le H\}$.
Its probability $p_H(r,k,a)$ is nonincreasing in the reserve share $r$, the seed $k$ and the loan rate $a$.
\end{lemma}

\begin{proposition}[Volume-maximising share]\label{prop:rhat}
Let supply be \eqref{eq:supply} with $A$ continuous and nonincreasing, $G$ continuous and strictly increasing, $V^d$ as in
Proposition~\ref{prop:eq}, and let the distribution of the loss rates $L_t$ have no atoms, so that $p_H$ is continuous. An
equilibrium then exists and is unique for every $(r,k)$.
\begin{enumerate}
\item A volume-maximising share $\hat r$ exists with $\hat r\,a^*(\hat r)\ge\EL$: at the maximum, the reserve's
  inflow is at least expected loss at the equilibrium rate.
\item If $\hat r$ is interior, $a^*(\hat r)\hat r>\EL$ and the primitives are differentiable, then
  \begin{equation}\label{eq:foc}
    s_y\cdot u\,a^*(\hat r)\;=\;s_p\cdot\pi_r,\qquad s_y=\frac{\partial\ln G(y-e)}{\partial y},\quad s_p=-\frac{A'(p)}{A(p)},\quad \pi_r=-\frac{\partial p_H}{\partial r},
  \end{equation}
  evaluated at the equilibrium: the yield that the marginal lender gives up for a point of share, valued at its
  hazard of leaving, equals the reduction in the loss probability, valued at the safety hazard.
\item Equilibrium volume $V^*$ is nondecreasing in the seed $k$ at every $r$.
\end{enumerate}
\end{proposition}

Part 1 says that a share whose inflow falls short of expected loss is never worth choosing: raising it costs
lenders nothing in the long run and makes the pool safer. Part 3 says that outside first-loss capital raises volume at every share.
Whether it should replace share depends on its cost, which Section~\ref{sec:numbers} compares.

% ---------------------------------------------------------------------------------------------
\section{Liquidity}\label{sec:liquidity}

Loans of term $T$ mature at a uniform rate if originations are staggered, so a fraction $dt/T$ of the loan book
repays in each interval $dt$. Repayments go first to the queue of withdrawals.

\begin{proposition}[Time to exit]\label{prop:liquidity}
Suppose lenders request a fraction $W\le1$ of deposits at once, loans repay on time and uniformly over the term, and
no new loan is made until the queue is paid. Let $q_0=(1-u)+c$ be the cash on hand per unit of deposits, with $c=R/D$ the
reserve's cash cushion. The run is paid in full after
\begin{equation}\label{eq:wait}
  w(W)=T\,\frac{(W-q_0)^+}{u}
\end{equation}
days, so a full run takes $T\,(1-c/u)$ days at most. A run of size $W$ is paid within $\bar w$ days if and only if
\begin{equation}\label{eq:ubound}
  u\le\frac{T\,(1-W+c)}{T-\bar w}.
\end{equation}
\end{proposition}

The pool is therefore liquid at the scale of one loan term: $\WaitFull$ days for a full run, $\WaitHalf$ days for a run
of half the deposits and $\WaitQuarter$ days for a quarter, at $u=\UBase\%$ and no cushion. A 10\% cushion
from the reserve cuts the full run to $\WaitFullCushion$ days. A tighter target bounds utilisation. A seven-day guarantee against
a run of half the deposits needs $u\le\UMaxHalfWeek\%$ ($\UMaxHalfWeekCushion\%$ with a 10\% cushion), against the contract's
90\% cap. By \eqref{eq:identity}, the lower utilisation raises the floor from \AFreeAtCap\% to \AFreeAtULiq\%: that is what the
guarantee costs borrowers (Table~\ref{tab:liquidity}). Late repayments lengthen the wait: if a fraction $\pi$ of maturing
loans does not repay on time, replace $u$ by $(1-\pi)u$ in \eqref{eq:wait}.

\begin{figure}[ht]
  \centering
  % Figure: days until a run of a given size is paid out, by utilisation (data/liquidity_curve.csv).
\begin{tikzpicture}
  \begin{axis}[
      width=0.82\linewidth, height=5.6cm,
      xmin=0, xmax=1, ymin=0, ymax=32,
      xlabel={run: share of deposits requested at once}, ylabel={days until paid in full},
      tick label style={font=\footnotesize}, label style={font=\small},
      legend style={font=\scriptsize, draw=none, fill=white, fill opacity=0.85, text opacity=1, at={(0.03,0.97)}, anchor=north west},
    ]
    \addplot[thick, red!70!black] table[x=w, y=u90, col sep=comma] {data/liquidity_curve.csv};
    \addlegendentry{$u=90\%$ (contract cap)}
    \addplot[thick, blue!65!black] table[x=w, y=u85, col sep=comma] {data/liquidity_curve.csv};
    \addlegendentry{$u=85\%$}
    \addplot[thick, green!45!black] table[x=w, y=u70, col sep=comma] {data/liquidity_curve.csv};
    \addlegendentry{$u=70\%$}
    \addplot[thick, blue!65!black, dashed] table[x=w, y=u85c, col sep=comma] {data/liquidity_curve.csv};
    \addlegendentry{$u=85\%$, reserve cushion $10\%$}
  \end{axis}
\end{tikzpicture}

  \caption{Days until a run is paid in full (Proposition~\ref{prop:liquidity}), $T=\TermDays$ days.}
  \label{fig:liquidity}
\end{figure}

\begin{table}[ht]
  \centering\small
  \caption{Utilisation allowed by a liquidity target, and the floor at that utilisation. For a run of $W$ of deposits
  and a wait of at most 3, 7 or 14 days: the largest utilisation (\%, capped at the contract's 90\%) and the loan
  rate at which lenders earn $e$ with $r\le r_c$ (\%).}
  \label{tab:liquidity}
  \begin{tabular}{rrrrrrr}
    \toprule
    & \multicolumn{2}{c}{3 days} & \multicolumn{2}{c}{7 days} & \multicolumn{2}{c}{14 days} \\
    \cmidrule(lr){2-3}\cmidrule(lr){4-5}\cmidrule(lr){6-7}
    $W$ (\%) & $u$ & floor & $u$ & floor & $u$ & floor \\
    \midrule
    \inputrows{data/table_liquidity.tex}
    \bottomrule
  \end{tabular}
\end{table}

% ---------------------------------------------------------------------------------------------
\section{What the reserve share buys borrowers}\label{sec:growth}

\begin{proposition}[Earned credit]\label{prop:growth}
A borrower with own credit $\ell$ who keeps its whole limit drawn on rolling loans earns dues at rate
$r\,a\,(\ell+d)$ and so holds a limit $\ell\,e^{r\,a\,t}$ after $t$ years. It doubles in $\ln 2/(r\,a)$ years, and each
dollar of earned credit costs $1/r$ dollars of interest.
\end{proposition}

At the production rate and $r=45\%$ the limit grows at $\GrowthAtFortyFive\%$ a year and doubles in $\DoubleAtFortyFive$ years; at
$65\%$ it grows at $\GrowthAtSixtyFive\%$ and doubles in $\DoubleAtSixtyFive$ years. Moving from 45\% to 65\% buys a borrower
\GrowthGain\ points of annual growth in its limit and costs lenders $\SlopePerPoint$ points of yield per point of share each year. Earned credit
is a slow engine of growth either way; faster growth has to come from the issuer's lines. We therefore treat the borrower's
benefit from the reserve as second-order and do not add it to the objective.

% ---------------------------------------------------------------------------------------------
\section{Numbers}\label{sec:numbers}

\subsection{Lenders' risk of losing principal}

Table~\ref{tab:safety} reports, for the inputs of Table~\ref{tab:inputs} and 200,000 simulated paths, the probability that
lenders' return is negative in some year of the first ten, with and without a seed. The identity of Lemma~\ref{lem:return}
holds on every path to within $\IdentityErr$. At $r=45\%$ the probability is \PTenAtFortyFive\%, against \PTenAtRC\% at $r_c$,
\PTenAtFifty\% at 50\% and \PTenAtSixtyFive\% at 65\%. The probability does not reach zero at any share, because a
bad year before the reserve has built up falls on lenders: the early-year protection depends mostly on the seed. A seed of 2\% of
deposits lowers the probability at 45\% to \PTenAtFortyFiveSeedTwo\%, and a seed of 5\% to \PTenAtFortyFiveSeedFive\%.

\begin{figure}[ht]
  \centering
  % Figure: probability that lenders lose principal in some year of ten, against the reserve share
% and the seed capital (data/safety_curve.csv).
\begin{tikzpicture}
  \begin{axis}[
      width=0.82\linewidth, height=5.8cm,
      xmin=0.2, xmax=0.7, ymin=0, ymax=16,
      xlabel={reserve share $r$}, ylabel={P(principal loss within 10 years) (\%)},
      tick label style={font=\footnotesize}, label style={font=\small},
      legend style={font=\scriptsize, draw=none, fill=white, fill opacity=0.85, text opacity=1, at={(0.97,0.97)}, anchor=north east},
    ]
    \addplot[thick, red!70!black] table[x=r, y=p_k0, col sep=comma] {data/safety_curve.csv};
    \addlegendentry{no seed}
    \addplot[thick, blue!65!black] table[x=r, y=p_k2, col sep=comma] {data/safety_curve.csv};
    \addlegendentry{seed $2\%$ of deposits}
    \addplot[thick, green!45!black] table[x=r, y=p_k5, col sep=comma] {data/safety_curve.csv};
    \addlegendentry{seed $5\%$}
    \addplot[dotted, gray] coordinates {(0.4176,0) (0.4176,16)};
    \node[font=\scriptsize, rotate=90, anchor=south] at (axis cs:0.4176,8) {$r_c$};
  \end{axis}
\end{tikzpicture}

  \caption{Probability that lenders lose principal in some year of ten, against the reserve share, by seed capital
  (200,000 paths).}
  \label{fig:safety}
\end{figure}

\begin{table}[ht]
  \centering\small
  \caption{Lender safety at the production rate. Columns: reserve share (\%); probability of a negative return in some
  year of ten (\%) with no seed, a 2\% seed and a 5\% seed; mean annual return over the ten years, no seed (\%);
  expected reserve after ten years, per unit of deposits (\%).}
  \label{tab:safety}
  \begin{tabular}{rrrrrr}
    \toprule
    $r$ & no seed & seed 2\% & seed 5\% & mean return & reserve \\
    \midrule
    \inputrows{data/table_safety.tex}
    \bottomrule
  \end{tabular}
\end{table}

Two features need comment. First, the long-run cost of a share up to $r_c$ is zero, but the ten-year mean return is not: it is
\MeanRC\% at $r_c$ against \MeanZero\% with no reserve, because the reserve is still filling and its balance, $\ResTenRC\%$ of
deposits after ten years, belongs to lenders who have not yet earned it. Early lenders pay for a reserve that later lenders enjoy. A seed
removes that transfer. Second, a mean--variance penalty does not justify a larger share. At $r=45\%$, a penalty with risk-aversion
coefficient 5, 20 or 100 is \MVPenFive, \MVPenTwenty\ or \MVPenHundred\ basis points of annual return, against \ShareCostFortyFive\ percentage points a year for
the move from $r_c$ to 45\% (Table~\ref{tab:seed}). Only a criterion tied to the loss of principal gives the share a role.

Table~\ref{tab:rsafe} gives the smallest share that keeps the probability at or below a tolerance $\alpha$. Combined with
Proposition~\ref{prop:rhat}, the share to choose is $\max\{r_c,\ r_{\mathrm{safe}}\}$.

\begin{table}[ht]
  \centering\small
  \caption{Smallest reserve share (\%) with a ten-year probability of losing principal at or below the tolerance;
  ``--'' where no share up to 70\% suffices. Shares below $r_c$ are free, so the share to choose is the larger of
  $r_c$ and the entry.}
  \label{tab:rsafe}
  \begin{tabular}{rrrr}
    \toprule
    tolerance (\%) & no seed & seed 2\% & seed 5\% \\
    \midrule
    \inputrows{data/table_rsafe.tex}
    \bottomrule
  \end{tabular}
\end{table}

\subsection{Seed capital against share}

Table~\ref{tab:seed} asks what seed, added at the share $r_c$ that costs lenders nothing, gives the protection of a higher share. A
share above $r_c$ costs lenders $u\,a\,(r-r_c)$ a year, per unit of deposits, for as long as the pool lives. A seed costs its
funder a return on the capital committed. Matching the protection of 45\% takes a seed of \SeedEquivFortyFive\% of deposits, which costs
\SeedCostFortyFive\% a year at a 15\% hurdle against $\ShareCostFortyFive\%$ for the share; matching 55\% takes a seed of
\SeedEquivFiftyFive\%, which costs \SeedCostFiftyFive\% against $\ShareCostFiftyFive\%$. At the funding rate as the cost of capital the seed is cheaper still. The seed does not
replace the reserve: it protects against the first years, while the reserve builds, and the model never releases it.

\begin{table}[ht]
  \centering\small
  \caption{Seed against share. Columns: the share (\%); its ten-year probability of losing principal (\%, no seed); the
  seed (\% of deposits) that gives the same probability at $r_c$; the yearly cost to lenders of the share above $r_c$ (\%);
  the yearly cost of the seed at a cost of capital equal to the funding rate and at 15\% (\%).}
  \label{tab:seed}
  \begin{tabular}{rrrrrr}
    \toprule
    $r$ & prob. & seed & share cost & seed at $e$ & seed at 15\% \\
    \midrule
    \inputrows{data/table_seed.tex}
    \bottomrule
  \end{tabular}
\end{table}

\subsection{Equilibrium}

Tables~\ref{tab:eqcurve} and~\ref{tab:eq} compute the equilibrium for the inputs of Table~\ref{tab:inputs}. Demand has constant elasticity $\eta$ and
lenders' margins are exponential with mean $\tau$, so that supply saturates as lenders' capacity is used. We
calibrate demand and lenders' capacity so that the production rate clears the market at $r_c$ with no seed, and report
volumes relative to that point: the levels are not predictions, only the response to $r$. The first
scenario has no safety requirement. Up to $r_c$ the share changes nothing, as Proposition~\ref{prop:identity} says; beyond it,
the equilibrium rate rises to \AFortyFiveYield\% at 45\% and to \ASixtyFiveYield\% at 65\%, and volume falls to
\VFortyFiveYield\ and \VSixtyFiveYield\ of its level at $r_c$. If the premium is held at \PremBps\ basis points, borrowers can
still ask for the same volume but lenders fund only \VFixedFortyFiveYield\ of it at 45\%, and nothing from about 60\% (where their
yield falls below the funding rate); Table~\ref{tab:eqcurve}.

With a safety requirement (tolerance exponential with mean $\bar\alpha=10\%$), more reserve brings in lenders who
would otherwise stay out, and volume peaks. In the base case the peak is at $\hat r=\RHatBase\%$, where volume is \VMaxBase\ of its
reference level, and volume stays within 1\% of that peak for shares between \LoNinetyNine\% and \HiNinetyNine\%. At 45\% it is
\VFortyFiveBase, at 65\% it is \VSixtyFiveBase. Across the safety scenarios of Table~\ref{tab:eq} the peak lies between \RHatMin\% and
\RHatMax\%, and volume falls steeply above about 55\%: the more elastic demand is, the faster. A seed leaves the peak
almost where it is and raises volume at every share (Proposition~\ref{prop:rhat}, part~3): at the peak by \SeedGainTwo\% for a seed of 2\% of deposits and by \SeedGainFive\% for a seed of 5\%.

\begin{figure}[ht]
  \centering
  % Figure: equilibrium loan volume and APR against the reserve share (data/equilibrium_curves.csv).
\begin{tikzpicture}
  \begin{groupplot}[
      group style={group size=2 by 1, horizontal sep=1.7cm},
      width=0.47\linewidth, height=5.4cm,
      xmin=0.3, xmax=0.7, xlabel={reserve share $r$},
      tick label style={font=\footnotesize}, label style={font=\small},
      legend style={font=\scriptsize, draw=none, fill=white, fill opacity=0.85, text opacity=1},
    ]
    \nextgroupplot[ylabel={equilibrium volume}, ymin=0.6, ymax=1.1, legend style={at={(0.03,0.03)}, anchor=south west}]
    \addplot[thick, black, dashed] table[x=r, y=v_yield, col sep=comma] {data/equilibrium_curves.csv};
    \addlegendentry{yield only}
    \addplot[thick, red!70!black] table[x=r, y=v_base, col sep=comma] {data/equilibrium_curves.csv};
    \addlegendentry{with safety, no seed}
    \addplot[thick, blue!65!black] table[x=r, y=v_base_k2, col sep=comma] {data/equilibrium_curves.csv};
    \addlegendentry{seed $2\%$}
    \addplot[thick, green!45!black] table[x=r, y=v_base_k5, col sep=comma] {data/equilibrium_curves.csv};
    \addlegendentry{seed $5\%$}
    \addplot[thick, red!70!black, dotted] table[x=r, y=v_h10, col sep=comma] {data/equilibrium_horizon_curves.csv};
    \addlegendentry{ten-year horizon, with safety}
    \nextgroupplot[ylabel={equilibrium APR (\%)}, ymin=11, ymax=18]
    \addplot[thick, black, dashed] table[x=r, y=a_yield, col sep=comma] {data/equilibrium_curves.csv};
    \addplot[thick, red!70!black] table[x=r, y=a_base, col sep=comma] {data/equilibrium_curves.csv};
    \addplot[thick, blue!65!black] table[x=r, y=a_base_k2, col sep=comma] {data/equilibrium_curves.csv};
    \addplot[thick, green!45!black] table[x=r, y=a_base_k5, col sep=comma] {data/equilibrium_curves.csv};
    \addplot[thick, red!70!black, dotted] table[x=r, y=a_h10, col sep=comma] {data/equilibrium_horizon_curves.csv};
  \end{groupplot}
\end{tikzpicture}

  \caption{Equilibrium volume (relative to the production volume at $r_c$) and APR against the reserve share; $\eta=1$,
  $\tau=1\%$, $\bar\alpha=10\%$.}
  \label{fig:eq}
\end{figure}

\begin{table}[ht]
  \centering\small
  \caption{Equilibrium by reserve share: loan APR (\%), volume relative to the production volume at $r_c$, and the volume
  lenders fund if the premium is held at \PremBps\ basis points (``fixed''). First, no safety requirement; then a safety
  requirement with $\bar\alpha=10\%$, no seed.}
  \label{tab:eqcurve}
  \begin{tabular}{rrrrrrr}
    \toprule
    & \multicolumn{3}{c}{yield only} & \multicolumn{3}{c}{with safety} \\
    \cmidrule(lr){2-4}\cmidrule(lr){5-7}
    $r$ & APR & volume & fixed & APR & volume & fixed \\
    \midrule
    \inputrows{data/table_eq_curve.tex}
    \bottomrule
  \end{tabular}
\end{table}

\begin{table}[ht]
  \centering\small
  \caption{The volume-maximising share by scenario. Columns: demand elasticity $\eta$; mean lenders' margin $\tau$ (\%); mean
  tolerated probability of losing principal $\bar\alpha$ (\%); seed (\% of deposits); volume-maximising share (\%); volume
  at the maximum; shares within 1\% of the maximum volume (\%); volume at 45\%, 55\% and 65\%. Volumes relative to
  the production volume at $r_c$ with no seed.}
  \label{tab:eq}
  \begin{tabular}{rrrrrrrrrr}
    \toprule
    $\eta$ & $\tau$ & $\bar\alpha$ & seed & $\hat r$ & max & within 1\% & 45\% & 55\% & 65\% \\
    \midrule
    \inputrows{data/table_equilibrium.tex}
    \bottomrule
  \end{tabular}
\end{table}

\subsection{Other books}

The best share depends on the book's default risk. Table~\ref{tab:robust} repeats the base case for annual default
probabilities of 3\% and 10\%, with the premiums of 600 and 1,400 basis points that the calibration of
\citet{pricingreserve2026} gives for them. The share that covers expected loss is 29.6\%, \RC\% and 57.6\%; the volume-maximising
share is about four to six points above it each time; and the shares within 1\% of the maximum do not overlap: 45\% is on the plateau for
the 5\% book, above it for the 3\% book, and below it for the 10\% book.

\begin{table}[ht]
  \centering\small
  \caption{Three books. Columns: annual default probability (\%); premium (bps); loan rate (\%); expected loss (\%);
  share that covers expected loss $r_c$ (\%); floor at no margin and $r\le r_c$ (\%); share at which lenders earn only the funding rate
  at the production rate (\%); volume-maximising share, base case (\%); shares within 1\% of the maximum (\%); volume at
  45\% and at 65\%, relative to the volume at $r_c$.}
  \label{tab:robust}
  \begin{tabular}{rrrrrrrrrrr}
    \toprule
    PD & prem. & APR & EL & $r_c$ & floor & $r$ for $e$ & $\hat r$ & within 1\% & 45\% & 65\% \\
    \midrule
    \inputrows{data/table_robust.tex}
    \bottomrule
  \end{tabular}
\end{table}

\begin{table}[ht]
  \centering\small
  \caption{Lenders' return and borrowers' earned credit at the production rate. Columns: reserve share $r$ (\%); lenders' long-run
  return~\eqref{eq:longrun} (\%); interest net of the reserve share, $u\,a(1-f-r)$ (\%); return above the funding rate (\%);
  growth rate of a borrower's limit, $r\,a$ (\%); years to double (Proposition~\ref{prop:growth}); dollars of interest per dollar of
  earned credit, $1/r$.}
  \label{tab:lender}
  \begin{tabular}{rrrrrrr}
    \toprule
    $r$ & return & net interest & above $e$ & growth & doubling & $1/r$ \\
    \midrule
    \inputrows{data/table_lender.tex}
    \bottomrule
  \end{tabular}
\end{table}

% ---------------------------------------------------------------------------------------------
\section{Limits and extensions}\label{sec:limits}

\begin{enumerate}
\item \emph{Elasticities and tolerances are scenarios.} The paper has no estimate of the pool's demand elasticity, of lenders'
  response to yield, or of the distribution of their tolerance for loss. The results that do not depend on them are the identity
  (Proposition~\ref{prop:identity}), the existence and monotonicity results, the incidence formula as a function of the
  elasticities, and part 1 of Proposition~\ref{prop:rhat}. The numerical plateau and its position depend on them. The
  pool's own flows can identify them: deposit inflows against the yield offered, and loan requests against the rate, once the pool has
  operated through some changes in both.
\item \emph{No data on default.} The 3, 5 and 10\% books are scenarios. The loss model is a one-factor Gaussian model with
  independent years, constant exposure, loss given default of 100\% and no seasoning \citep{vasicek2002,pricingreserve2026}.
\item \emph{Higher rates may raise default.} If an increase of one point in the loan rate raises the annual default probability by $\gamma$
  points, lenders' long-run yield in the regime $\EL(a)>a\,r$ bends back when $\gamma>(1-f)/\kappa=\GammaBackward$, where $\kappa=\KappaLoss$ is the
  slope of expected loss in the default probability, and credit is rationed \citep{stiglitz1981}. While the reserve covers
  losses ($a\,r>\EL(a)$) lenders' yield $u\,a\,(1-f-r)$ is unaffected, but the lock fails once $\kappa\gamma>r$: $\gamma>\GammaProtectFortyFive$ at $r=45\%$ and $\gamma>\GammaProtectSixtyFive$ at 65\%. Neither value has been measured.
\item \emph{Static, steady-state analysis.} The model compares equilibria. It does not model the path of deposits and loans, entry of
  competing pools, or the response of other protocols' rates.
\item \emph{Safety-first with one criterion.} The probability of a negative return in a year of the first ten is one of several
  reasonable criteria. It ignores the size of the loss and treats a loss of one cent as one of the whole deposit.
\item \emph{Liquidity assumes staggered maturities and on-time repayment.} Synchronised maturities change the shape of \eqref{eq:wait}
  but not its bound for a full run; late repayment lengthens it.
\item \emph{The reserve's cash is idle.} The model leaves the cash that the reserve holds outside the lending base, as the contract does. Lending
  it would raise interest and risk, and the analysis here does not cover it.
\item \emph{Synthetic accounts.} Nothing here is evidence about the welfare of borrowers.
\end{enumerate}

% ---------------------------------------------------------------------------------------------
\section{Implications for the choice of the share}\label{sec:implications}

The analysis changes how the question is posed, and in this scenario supports six statements.
\begin{enumerate}
\item \emph{The best share is a plateau near the share that covers expected loss, not a point.} For the 5\% book it extends from
  \LoNinetyNine\% to \HiNinetyNine\%, with its peak about four to six points above $r_c$ in each of three books. The share that covers
  expected loss is itself a function of the book: \RC\% at 5\%, 29.6\% at 3\%, 57.6\% at 10\%.
\item \emph{Beyond the plateau the cost rises quickly.} Each point above $r_c$ costs lenders $\SlopePerPoint$ points of yield per
  year; the equilibrium rate must rise, and volume falls by about \VolPerPointFifty\% per point near 50\% in the yield-only scenario (formula~\eqref{eq:incidence}) and faster above. Above $\RMaxLenders\%$
  lenders earn less than the funding rate at the production rate.
\item \emph{Seed capital is cheaper protection than share.} It protects the first years, when lenders bear most of the risk
  and reserve funded by interest has not yet built up.
\item \emph{The premium and the share move together.} If the share rises and the premium does not, volume falls to what lenders will
  fund (Table~\ref{tab:eqcurve}, ``fixed''). The incidence formula gives the premium that restores volume.
\item \emph{The share is a weak lever for growing borrowers' credit.} Dues add a few percent a year to a limit (Table~\ref{tab:lender}).
\item \emph{The pool is liquid at the scale of one loan term.} A stricter promise to lenders bounds utilisation and raises the rate floor.
\end{enumerate}
These are results of the model and its scenarios. The choice of a share in production should also weigh what the model leaves out: the
protocol's governance, the risk of the issuer's lines, the cost and availability of seed capital, and evidence from the pool itself.

\section{Conclusion}

A reserve share sets a wedge between what borrowers pay and what lenders receive, and a reserve that is never released makes the wedge
permanent once it exceeds expected loss. We have shown how the wedge divides between the two sides of the market, why a share below
expected loss is free, why a share a few points above it is worth having when lenders care about losing principal, and why
a larger share reduces lending. The pool's liquidity follows from the term of its loans. What remains is measurement: the pool's own flows
can pin down the elasticities on which the position of the plateau depends.

\paragraph{Disclosure of AI authorship} This paper was written by Claude Code, an AI agent made by Anthropic, working with the project's
human operator, who set the research direction and asked for a demand and equilibrium model of the reserve share. The simulation of lenders'
returns reuses the loss model of the calibration code written by AI agents working with the project; the equilibrium, safety and liquidity
computations are this paper's own. All of it is public and reproducible.

{\sloppy\paragraph{Data and code availability} The loss model and its calibration, at commit b725a85, are in
\texttt{analysis/credit\_risk} of \url{https://github.com/scottonchain/microcredit-contract}. The script that produces every table, curve and
number of this paper, with the commands to reproduce them, is in \texttt{lending-equilibrium/scripts} of
\url{https://github.com/scottonchain/microcredit-theory}.\par}

\bibliographystyle{elsarticle-harv}
\bibliography{references}

\appendix
\section{Proofs}

\paragraph{Proof of Lemma~\ref{lem:return}}
If $R_{t-1}+u\,a\,r\ge u\,L_t$, then $R_t=R_{t-1}+u\,a\,r-u\,L_t$ and $\ell_t=0$. Otherwise $R_t=0$ and
$\ell_t=u\,L_t-R_{t-1}-u\,a\,r$, so again $R_t-R_{t-1}=u\,a\,r-u\,L_t+\ell_t$. Hence
$Y_t=u\,a(1-f-r)-\ell_t=u\,a(1-f)-u\,L_t-(R_t-R_{t-1})$, and summing gives the identity.

Let $\xi_t=u\,a\,r-u\,L_t$, $S_t=\sum_{s\le t}\xi_s$ and $\mu=u(a\,r-\EL)$. Unrolling the recursion, $R_T=\max\{k+S_T,\ \max_{1\le j\le T}(S_T-S_j)\}$.
The $\xi_t$ are independent and bounded, so $\sum_t\mathrm{Var}(\xi_t)/t^2<\infty$ and Kolmogorov's strong law gives $S_n/n\to\mu$
almost surely \citep{billingsley1995}. If $\mu>0$, $\min_jS_j$ is finite, $R_T=S_T+O(1)$ and $R_T/T\to\mu$. If $\mu\le0$, then for any $\epsilon>0$ almost surely $|S_n-n\mu|\le\epsilon n$ for large $n$, so
$S_T-S_j\le(T-j)\mu+\epsilon(T+j)\le2\epsilon T$ and $0\le R_T\le\max\{k+S_T,2\epsilon T\}$, whence $R_T/T\to0$. In both cases $R_T/T\to\max\{\mu,0\}$. The loss rates also obey the strong law, $T^{-1}\sum L_t\to\EL$. Dividing the identity by $T$,
$T^{-1}\sum Y_t\to u\,a(1-f)-u\,\EL-\max\{u\,a\,r-u\,\EL,0\}=y(a,r)$. \hfill$\square$

\paragraph{Proof of Proposition~\ref{prop:identity}}
$y(a,r)=\min\{u[a(1-f)-\EL],\ u\,a\,(1-f-r)\}$ is the minimum of two strictly increasing continuous functions of $a$ (as $r<1-f$), so it is strictly increasing, and its inverse at $e+x$ is the maximum of the two inverses, which is \eqref{eq:floor}. Identity \eqref{eq:identity} is $y=e+x$ rearranged. The first branch of \eqref{eq:floor} does not depend on $r$ and is the larger one exactly when $a_0(x)\,r\le\EL$, that is, for $r\le\bar r(x)$. The second is strictly increasing and convex in $r$, with derivative $(e+x)/(u(1-f-r)^2)=\underline a/(1-f-r)$, and unbounded as $r\to1-f$. \hfill$\square$

\paragraph{Proof of Proposition~\ref{prop:eq}}
Let $\Phi(a)=y(a,r)-e-\psi(V^d(a))$, with $\Phi=-\infty$ where $V^d(a)\ge\bar V$. The function $y(\cdot,r)$ is continuous and strictly increasing. Since $V^d$ is decreasing and $\psi$ increasing, $\psi(V^d(a))$ is nonincreasing, so $\Phi$ is strictly increasing. At $a=0$, $y=-u\,\EL<e$, so $\Phi<0$. As $a\to\infty$, $y\to\infty$ while $\psi(V^d(a))$ is bounded above by its value at any finite $a$, so $\Phi>0$ for large $a$. By continuity $\Phi$ has a root, unique because $\Phi$ is strictly increasing; at it, $y=e+\psi(V^*)$ with $V^*=V^d(a^*)$, which is \eqref{eq:eqrate} after rearranging as in Proposition~\ref{prop:identity}. For $r'>r$, $y(a,r')\le y(a,r)$, so $\Phi(a^*(r);r')\le0$ and $a^*(r')\ge a^*(r)$; the inequality is strict where $a^*(r)\,r>\EL$ because $y$ then strictly decreases in $r$. Then $V^*=V^d(a^*)$ and $\psi(V^*)$ are nonincreasing. \hfill$\square$

\paragraph{Proof of Proposition~\ref{prop:incidence}}
In the regime $a\,r>\EL$, $y=u\,a\,(1-f-r)$, so $d\ln y=d\ln a-dr/(1-f-r)$. Along the equilibrium, $d\ln V^s=\varepsilon\,d\ln y$ and $d\ln V^d=-\eta\,d\ln a$ are equal. Hence $(\varepsilon+\eta)\,d\ln a=\varepsilon\,dr/(1-f-r)$, which gives the first formula. Then $d\ln V^*=-\eta\,d\ln a$ and $d\ln y^*=d\ln a-dr/(1-f-r)=-\eta/(\eta+\varepsilon)\cdot dr/(1-f-r)$. \hfill$\square$

\paragraph{Proof of Lemma~\ref{lem:mono}}
If $\ell_t>0$ then $R_t=0$ and $Y_t=u\,a(1-f)-u\,L_t+R_{t-1}$; if $\ell_t=0$ then $Y_t=u\,a(1-f-r)>0$. So $Y_t<0$ if and only if $u\,L_t>u\,a(1-f)+R_{t-1}$. By induction on \eqref{eq:recursion}, on every path each $R_t$ is nondecreasing in $r$, $k$ and $a$. The right side of the inequality is nondecreasing in all three, so the event shrinks as any of them rises, and so does its probability. \hfill$\square$

\paragraph{Proof of Proposition~\ref{prop:rhat}}
(1) Equilibrium solves $\Phi(a;r)=\ln V^s(y(a,r),p_H(r,k,a))-\ln V^d(a)=0$. Existence of an equilibrium and its uniqueness follow as in Proposition~\ref{prop:eq}: $\Phi$ is strictly increasing in $a$ (by Lemma~\ref{lem:mono}, $p_H$ falls as $a$ rises, so supply rises), equals $-\infty$ for $y\le e$ and is positive for large $a$. The root $a^*(r)$ is continuous in $r$, so $V^*$ is continuous; and since $(1-f-r)\,a^*\ge e/u$ by \eqref{eq:identity}, $a^*\to\infty$ and $V^*\to0$ as $r\to1-f$, so $V^*$ attains its maximum. Take a maximiser $r_0$. If $a^*(r_0)\,r_0\ge\EL$ we are done. Otherwise let $r_1=\EL/a^*(r_0)>r_0$ and $a=a^*(r_0)$. Then $\max\{\EL,a\,r_1\}=\EL=\max\{\EL,a\,r_0\}$, so the yield is the same, and $p_H(r_1,k,a)\le p_H(r_0,k,a)$ by Lemma~\ref{lem:mono}, so supply at $a$ is no smaller. As supply increases in $a$ (through $y$ and $p_H$) and demand decreases, the equilibrium rate at $r_1$ is no larger, so $V^*(r_1)\ge V^*(r_0)$ and $r_1$ is also a maximiser, with $a^*(r_1)\,r_1\ge\EL$.
(2) At an interior maximiser $dV^*/dr=0$, hence $da^*/dr=0$, so $\partial\Phi/\partial r=0$: $\partial\ln V^s/\partial r=-s_y\,u\,a+s_p\,\pi_r=0$, using $\partial y/\partial r=-u\,a$ in the regime $a\,r>\EL$, which is \eqref{eq:foc}.
(3) By Lemma~\ref{lem:mono}, $p_H$ is nonincreasing in $k$, so supply is nondecreasing in $k$ at every $(a,r)$, and the equilibrium rate is nonincreasing and volume nondecreasing. \hfill$\square$

\paragraph{Proof of Proposition~\ref{prop:liquidity}}
The pool's cash is $(1-u)D+R$ and lenders' claims total $D$. A run of $W\,D$ takes the cash first, leaving $(W-q_0)^+D$ to be paid from repayments. Loans of total $u\,D$ mature uniformly over $T$ days, so repayments arrive at $u\,D/T$ a day, all to the queue. The time to pay $(W-q_0)^+D$ is $T(W-q_0)^+/u$. For $W=1$ this is $T(u-c)/u=T(1-c/u)$. Setting $w(W)\le\bar w$ and solving for $u$ gives \eqref{eq:ubound}. \hfill$\square$

\paragraph{Proof of Proposition~\ref{prop:growth}}
Interest paid on a drawn limit $\ell+d$ is $a(\ell+d)$ a year, of which the share $r$ is added to $d$. So $\dot d=r\,a\,(\ell+d)$ and $\ell+d=\ell\,e^{r\,a\,t}$. The limit doubles when $r\,a\,t=\ln2$, and one dollar of dues takes $1/r$ dollars of interest. \hfill$\square$

\end{document}
